\documentclass[a4paper]{article}
\usepackage{amsfonts}
\usepackage{amsmath}
\usepackage{amssymb}
\usepackage{graphicx}     

\begin{document}
  
\noindent \large Auteur: Thomas Eertmans \\
\noindent \large Studierichting: Informatica\\
\noindent \large Oefeningenreeks 3F nr. 1.3\\

\medskip

\normalsize

$f(x) = \dfrac{x^3+2x^2+2x-2}{x^2+1}$ \\

$D = 0^2 - 4 \cdot 1 \cdot 1 < 0$ , dus geen VA. \\

HA: $\lim\limits_{x\rightarrow\infty} \dfrac{x^3+2x^2+2x-2}{x^2+1} = \dfrac{\infty}{\infty}$ \\

$\overset{\mathrm{H}}{=} \lim\limits_{x\rightarrow\infty} \dfrac{3x^2+4x+2}{2x} = \dfrac{\infty}{\infty}$ \\

$\overset{\mathrm{H}}{=} \lim\limits_{x\rightarrow\infty} \dfrac{6x+4}{2} = \infty$ , dus geen HA. \\

SA: $\lim\limits_{x\rightarrow\infty} \dfrac{\dfrac{x^3+2x^2+2x-2}{x^2+1}}{x} = \lim\limits_{x\rightarrow\infty} \dfrac{x^3+2x^2+2x-2}{x^3+x} = \dfrac{\infty}{\infty}$ \\

$\overset{\mathrm{H}}{=} \lim\limits_{x\rightarrow\infty} \dfrac{3x^2+4x+2}{3x^2+1} = \dfrac{\infty}{\infty}$ \\

$\overset{\mathrm{H}}{=} \lim\limits_{x\rightarrow\infty} \dfrac{6x+4}{6x} = \dfrac{\infty}{\infty}$ \\

$\overset{\mathrm{H}}{=} \lim\limits_{x\rightarrow\infty} \dfrac{6}{6} = 1 = a$ \\

$b = \lim\limits_{x\rightarrow\infty} \dfrac{x^3+2x^2+2x-2}{x^2+1} - x$ \\

$= \lim\limits_{x\rightarrow\infty} \dfrac{x^3+2x^2+2x-2-x^3-x}{x^2+1}$ \\

$= \lim\limits_{x\rightarrow\infty} \dfrac{2x^2+x-2}{x^2+1} = \dfrac{\infty}{\infty}$ \\

$\overset{\mathrm{H}}{=} \lim\limits_{x\rightarrow\infty} \dfrac{4x+1}{2x} = \dfrac{\infty}{\infty}$ \\

$\overset{\mathrm{H}}{=} \lim\limits_{x\rightarrow\infty} \dfrac{4}{2} = 2$ \\

$SA = x+2$ \\

Ligging tegenover de schuine asymptoot: \\

$\dfrac{x^3+2x^2+2x-2}{x^2+1} - x+2$ \\

$= \dfrac{x^3+2x^2+2x-2-(x+2)(x^2+1)}{x^2+1}$ \\

$= \dfrac{x^3+2x^2+2x-2-x^3-x-2x^2-2}{x^2+1} = \dfrac{x-4}{x^2+1}$

$\begin{array}{r | c c c c c}
x & -\infty &  & 4 & & + \infty \\ \hline
\dfrac{x-4}{x^2+1} & & - & 0 & + \\
\end{array}$ \\


\begin{thebibliography}{999}
%\bibitem{a} Referentie
\end{thebibliography}
\end{document} 
